High-performance computing (HPC) data management systems, such as parallel file systems, expose read/write interfaces rather than orchestrating data movement across a workflow. Staging, converting, and reducing data is left to users, and that burden grows with the scale and complexity of the workflow.

We present Data Flyway, a data management runtime that performs \emph{in-flight processing}: transformation and analysis operations applied to data while it moves from source to destination. Data Flyway represents a processing pipeline as a directed graph that is registered with an I/O resource. Any subsequent read or write on that resource traverses the graph, so data arrives at its destination already in the desired form, or accompanied by derived results, without the application invoking anything explicitly. Because a pipeline is a graph rather than a fixed chain, the runtime selects a path at execution time based on current resource availability and data movement cost.

We implement Data Flyway in Proactive Data Containers (PDC), a runtime data management library, and evaluate it on vector-magnitude analysis scaled to 128 nodes (4,096 MPI ranks, 4.1\,TB total) and curl-vorticity-magnitude analysis. At 4,096 ranks, Data Flyway reduces end-to-end time by 81\% relative to ADIOS2 and by 94\% relative to an HDF5 posthoc pipeline.